\chapter{Requirements Analysis}
\label{ch:requirements}

This chapter turns the scenario introduced in \cref{ch:introduction} into a
concrete set of requirements. \Cref{sec:stakeholders} identifies the
stakeholders and actors of the system. \Cref{sec:usecases} presents the
core use cases. \Cref{sec:functional-req} lists the functional
requirements. \Cref{sec:nonfunctional-req} lists the non-functional
requirements, and \cref{sec:offline-req} isolates, as a separate section,
the requirements that are specific to offline operation -- since these are
the requirements that most distinguish CampusFix from a conventional CRUD
application and that drive most of the design decisions in
\cref{ch:design}. \Cref{sec:data-req} specifies the data requirements for a
report record, and \cref{sec:constraints} states the constraints and
assumptions under which the requirements were elicited.

\section{Stakeholders and Actors}
\label{sec:stakeholders}

\begin{table}[htbp]
\centering
\caption{Stakeholders and actors of the CampusFix system.}
\label{tab:stakeholders}
\begin{tabularx}{\textwidth}{@{}l X@{}}
\toprule
\textbf{Actor} & \textbf{Description} \\
\midrule
\textbf{Reporter} & Any authenticated campus user (student, faculty, staff)
who creates, views, and manages maintenance reports through the mobile
app. This is the primary actor and the only actor implemented in the
current system. \\
\textbf{Facilities staff} & Consumer of the reports on the backend side.
Out of scope for the mobile application itself, but constrains the shape
of the data stored in Firestore (\cref{sec:data-req}), since that data must
eventually be actionable by this actor through a separate
dashboard/back-office, which is not part of this thesis. \\
\textbf{Synchronization module} & A non-human actor internal to the
system: the background component that moves data between the local
database and Firestore. Modeled as an actor because it initiates actions
(uploads, retries) independently of direct user input. \\
\textbf{Connectivity monitor} & A non-human actor that observes the
device's network state and notifies the synchronization module of
changes. \\
\bottomrule
\end{tabularx}
\end{table}

\section{Use Case Overview}
\label{sec:usecases}

\Cref{fig:usecase} summarizes the use cases supported by the Reporter
actor; \cref{tab:usecase-detail} details the three use cases most relevant
to the offline-first behaviour of the system.

\begin{figure}[htbp]
  \centering
  \begin{tikzpicture}[
      usecase/.style={draw, ellipse, minimum width=3.4cm, minimum height=1cm, align=center, font=\small},
      actor/.style={font=\small},
      >=Latex
    ]
    % Actor
    \node[actor] (actor) at (0,0) {
      \begin{tikzpicture}
        \draw (0,0.5) circle (0.2);
        \draw (0,0.3) -- (0,-0.4);
        \draw (-0.3,-0.1) -- (0.3,-0.1);
        \draw (0,-0.4) -- (-0.25,-0.9);
        \draw (0,-0.4) -- (0.25,-0.9);
      \end{tikzpicture}\\ Reporter
    };

    \node[usecase] (login) at (4,3.2) {Log in / Sign up};
    \node[usecase] (create) at (4,2.0) {Create report};
    \node[usecase] (view) at (4,0.8) {View reports\\(submitted \& pending)};
    \node[usecase] (retry) at (4,-0.4) {Retry failed sync};
    \node[usecase] (status) at (4,-1.6) {Check sync status};
    \node[usecase] (offline) at (8.4,0.8) {Work offline};

    \draw (actor) -- (login);
    \draw (actor) -- (create);
    \draw (actor) -- (view);
    \draw (actor) -- (retry);
    \draw (actor) -- (status);
    \draw[dashed,->] (create) -- node[above,font=\scriptsize]{\guillemotleft include\guillemotright} (offline);
    \draw[dashed,->] (view) -- node[below,font=\scriptsize]{\guillemotleft include\guillemotright} (offline);
  \end{tikzpicture}
  \caption{Use case diagram for the Reporter actor. "Work offline" is
  included by every core use case: none of them require connectivity to
  complete.}
  \label{fig:usecase}
\end{figure}

\begin{table}[htbp]
\centering
\caption{Detail of the three use cases most specific to offline-first
behaviour.}
\label{tab:usecase-detail}
\small
\begin{tabularx}{\textwidth}{@{}l X@{}}
\toprule
\textbf{Create report} & \textbf{Description} \\
\midrule
Actor & Reporter \\
Precondition & User is logged in (a session may be cached locally; see
\cref{sec:offline-req}). \\
Main flow & User opens "New report," fills in title, category,
description, optionally attaches a photo and a location, and submits. \\
Postcondition & A row exists in the local database with
\texttt{sync\_status = pending}. The report is immediately visible in the
report list, tagged as pending. \\
Offline behaviour & Identical to the online flow; no step in the main flow
requires network access. \\
\midrule
\textbf{View submitted and pending reports} & \textbf{Description} \\
\midrule
Actor & Reporter \\
Main flow & User opens the report list; the app displays all reports
belonging to the user, each annotated with its current sync status. \\
Offline behaviour & The list is read entirely from the local database; it
is always available and always reflects the true local state, including
reports not yet on the server. \\
\midrule
\textbf{Retry failed synchronization} & \textbf{Description} \\
\midrule
Actor & Reporter \\
Precondition & At least one report has \texttt{sync\_status = failed}. \\
Main flow & User triggers a manual sync (e.g. pull-to-refresh or a "retry"
action); the app re-attempts upload of all pending/failed reports. \\
Postcondition & Each report's status is updated to \texttt{synced} or
remains \texttt{failed}, depending on the outcome. \\
\bottomrule
\end{tabularx}
\end{table}

\section{Functional Requirements}
\label{sec:functional-req}

\Cref{tab:functional-req} lists the functional requirements (FR), grouped
by feature area and prioritized using the MoSCoW scale (Must, Should,
Could, Won't -- for this iteration).

\begin{longtable}{@{}p{1.4cm}p{8.6cm}p{1.6cm}@{}}
\caption{Functional requirements.}
\label{tab:functional-req}\\
\toprule
\textbf{ID} & \textbf{Requirement} & \textbf{Priority} \\
\midrule
\endfirsthead
\toprule
\textbf{ID} & \textbf{Requirement} & \textbf{Priority} \\
\midrule
\endhead
FR-1 & The system shall let a user create an account and log in with an
email and password. & Must \\
FR-2 & The system shall let a logged-in user create a report consisting of
a title, a category, a free-text description, an optional photo, a
location (building/room), and an automatically-assigned creation
timestamp. & Must \\
FR-3 & The system shall persist every report locally as soon as it is
created, independent of network state. & Must \\
FR-4 & The system shall let a user view the list of their own reports,
both those already synchronized and those still pending. & Must \\
FR-5 & The system shall let a user open a single report to see its full
detail, including its current synchronization status. & Must \\
FR-6 & The system shall let a user edit a previously created report (own
reports only). & Should \\
FR-7 & The system shall let a user delete a previously created report. &
Should \\
FR-8 & The system shall detect when network connectivity becomes
available and automatically attempt to synchronize pending reports. &
Must \\
FR-9 & The system shall let a user manually trigger synchronization. &
Must \\
FR-10 & The system shall display, for every report, one of the following
synchronization statuses: \texttt{pending}, \texttt{synced},
\texttt{failed}, \texttt{deleted\_pending}. & Must \\
FR-11 & The system shall retry synchronization of \texttt{failed} reports
without requiring the user to re-enter data. & Must \\
FR-12 & The system shall apply a defined conflict-resolution policy
(last-write-wins, \cref{sec:conflict-design}) when a report has been
modified both locally and on the server. & Should \\
FR-13 & The system shall upload photos attached to a report to the
backend's file storage and associate the resulting URL with the report
once online. & Could \\
FR-14 & The system shall support attaching a location either via device
GPS or via manual building/room selection. & Could \\
\bottomrule
\end{longtable}

\section{Non-Functional Requirements}
\label{sec:nonfunctional-req}

Non-functional requirements are grouped following the quality
characteristics of ISO/IEC 25010~\cite{iso25010}: functional suitability
is covered by \cref{sec:functional-req}, so this section focuses on
performance efficiency, reliability, usability, and portability, which are
the characteristics most affected by the offline-first constraint.

\begin{longtable}{@{}p{1.4cm}p{5.4cm}p{5.0cm}@{}}
\caption{Non-functional requirements.}
\label{tab:nonfunctional-req}\\
\toprule
\textbf{ID} & \textbf{Requirement} & \textbf{Rationale / metric} \\
\midrule
\endfirsthead
\toprule
\textbf{ID} & \textbf{Requirement} & \textbf{Rationale / metric} \\
\midrule
\endhead
NFR-1 & Creating a report locally shall complete in under 200\,ms,
perceived as instantaneous by the user. & Must not depend on network
round-trip time; measured as the time from tapping "Save" to the report
appearing in the list. \\
NFR-2 & The application shall remain fully navigable and operable with the
device in airplane mode. & Verified by the "offline creation" scenario in
\cref{ch:evaluation}. \\
NFR-3 & Synchronization of a single report shall not block the UI thread.
& Implemented via coroutines running off the main dispatcher
(\cref{ch:implementation}). \\
NFR-4 & The application shall not lose a locally-created report even if
the app is killed before synchronization completes. & Verified by
persisting to disk (Room/SQLite) on creation, not to in-memory state. \\
NFR-5 & The synchronization status of a report shall be visible within
one interaction of viewing the report (list or detail). & Usability
requirement, evaluated in \cref{sec:usability-eval}. \\
NFR-6 & The application shall share a single codebase for the Android and
iOS clients for all business logic, UI, and persistence code, with
platform-specific code limited to the thin platform integration layer. &
Maintainability; measured qualitatively in \cref{ch:implementation} by the
proportion of code in \texttt{commonMain}. \\
NFR-7 & The synchronization protocol shall be idempotent: re-running
synchronization after a partial failure shall not create duplicate remote
records. & Achieved via deterministic document IDs
(\cref{sec:sync-design}). \\
NFR-8 & Authentication credentials/session shall be cached locally so
that a returning user does not need connectivity to reach the
authenticated part of the app. & See \cref{sec:offline-req}, OFR-2. \\
\bottomrule
\end{longtable}

\section{Offline-Specific Requirements}
\label{sec:offline-req}

The requirements in this section are the ones that make CampusFix an
offline-first application rather than an online application with a cache.
They are listed separately, following the recommendation in
\cref{sec:offline-first-background}, precisely because they are easy to
omit from a conventional requirements document and expensive to retrofit.

\begin{longtable}{@{}p{1.6cm}p{9.0cm}@{}}
\caption{Offline-specific requirements.}
\label{tab:offline-req}\\
\toprule
\textbf{ID} & \textbf{Requirement} \\
\midrule
\endfirsthead
\toprule
\textbf{ID} & \textbf{Requirement} \\
\midrule
\endhead
OFR-1 & \textbf{Local-first writes.} Every write operation (create,
update, delete of a report) shall be applied to the local database first
and shall succeed regardless of network state; the network is never on
the critical path of a write. \\
OFR-2 & \textbf{Offline session continuity.} A user who has previously
logged in while online shall remain able to use the app (i.e. reach the
authenticated screens) while offline, using a cached session, without
being forced to re-authenticate against the network. \\
OFR-3 & \textbf{Deterministic local identity.} Every report shall receive
a locally-generated, globally unique identifier (\texttt{local\_id}) at
creation time, independent of any server-assigned identifier, so that the
report can be referenced, displayed, and edited before it has ever
reached the server. \\
OFR-4 & \textbf{Explicit sync state.} Every report shall carry an explicit
\texttt{sync\_status} field, persisted locally, distinguishing at least
\texttt{pending} (not yet uploaded), \texttt{synced} (confirmed on
server), \texttt{failed} (upload attempted and rejected/errored), and
\texttt{deleted\_pending} (deleted locally, deletion not yet propagated).
\\
OFR-5 & \textbf{Connectivity-triggered synchronization.} The
synchronization module shall subscribe to connectivity change events and
shall automatically attempt to synchronize all pending work items when
the device transitions from offline to online. \\
OFR-6 & \textbf{Bounded, resumable sync.} Synchronization of $N$ pending
reports shall be resumable: if interrupted after $k<N$ reports have been
uploaded, a subsequent sync attempt shall upload only the remaining
$N-k$, not re-upload the first $k$. \\
OFR-7 & \textbf{Conflict resolution policy.} When the synchronization
module detects that a report has diverging local and remote versions, it
shall apply a defined, documented policy (last-write-wins on
\texttt{updated\_at}, \cref{sec:conflict-design}) rather than an undefined
or silently inconsistent outcome. \\
OFR-8 & \textbf{No data loss on failure.} A failed synchronization attempt
shall never delete or corrupt the local copy of a report; the report shall
remain \texttt{pending}/\texttt{failed} and retriable. \\
\bottomrule
\end{longtable}

\section{Data Requirements}
\label{sec:data-req}

\Cref{tab:report-fields} specifies the fields required on a report record,
consolidating the "Core features" and "Synchronization logic" sections of
the project brief. This table is the basis for the local schema and remote
document shape designed in \cref{sec:local-storage-design}.

\begin{longtable}{@{}p{2.6cm}p{2.4cm}p{6.6cm}@{}}
\caption{Required fields of a report record.}
\label{tab:report-fields}\\
\toprule
\textbf{Field} & \textbf{Type} & \textbf{Notes} \\
\midrule
\endfirsthead
\toprule
\textbf{Field} & \textbf{Type} & \textbf{Notes} \\
\midrule
\endhead
\texttt{id} & integer & Local auto-increment primary key; never leaves the
device. \\
\texttt{local\_id} & UUID string & Client-generated identity, stable across
the report's lifetime; doubles as the remote document ID when no
server-assigned ID exists yet. \\
\texttt{remote\_id} & string, nullable & Firestore document ID once
synchronized; \texttt{null} until first successful upload. \\
\texttt{title} & string & Required, non-empty. \\
\texttt{category} & enum/string & One of a fixed set (projector, furniture,
network, HVAC, lab equipment, other). \\
\texttt{description} & string & Free text, required. \\
\texttt{location} & string & Building/room, free text or picker-selected. \\
\texttt{photo\_path} & string, nullable & Absolute local file path to an
attached photo. \\
\texttt{remote\_photo\_url} & string, nullable & Download URL once the
photo has been uploaded to remote storage. \\
\texttt{owner} & string & Email of the creating user; used for
per-user filtering and, informally, as the record's ownership token. \\
\texttt{created\_at} & epoch millis & Set once, at creation. \\
\texttt{updated\_at} & epoch millis & Updated on every local or remote
write; drives the last-write-wins policy. \\
\texttt{sync\_status} & enum & \texttt{pending} $\mid$ \texttt{synced}
$\mid$ \texttt{failed} $\mid$ \texttt{deleted\_pending}. \\
\bottomrule
\end{longtable}

\section{Constraints and Assumptions}
\label{sec:constraints}

\begin{itemize}
  \item \textbf{Technology stack.} The client is constrained to Kotlin
  Multiplatform targeting Android and iOS from a shared codebase; the
  backend is constrained to Firebase (Firestore for data, Firebase
  Authentication for identity), as specified in the project brief. These
  are treated as hard constraints rather than choices evaluated against
  alternatives.
  \item \textbf{Single-owner records.} A report is assumed to have exactly
  one meaningful "owner" at a time (the creating user); multi-user
  concurrent editing of the same report by two \emph{different} users is
  out of scope, though the same-user, two-device case is in scope for
  conflict handling (\cref{sec:conflict-design}).
  \item \textbf{Trusted client.} The system assumes an honest client: the
  \texttt{owner} field is set by the client and used for filtering, not as
  a security boundary. Backend authorization rules (Firestore security
  rules) are discussed as a limitation in \cref{ch:conclusions} rather than
  implemented in full.
  \item \textbf{Scale.} The system is designed and evaluated for the scale
  of a single university campus (hundreds to low thousands of reports),
  not for a nationwide, high-write-throughput deployment.
  \item \textbf{No offline maps/geocoding.} Location capture relies on
  device GPS coordinates and/or a static, bundled list of buildings; it
  does not depend on an online maps or geocoding service, so it does not
  introduce an additional offline dependency.
\end{itemize}
